\chapter{Preliminaries}

\section{Group theory}

In the following section we will introduce some group theoretic notions that
be used throughout this thesis. We assume that the Reader is familiar with
the notion of a group.

% This is from Dummit & Foote's Algebra, page 167
\begin{definition}[Dual of a group]
    Let \( G \) be a group. The \emph{dual group} \( \hat{G} \) is the
    set of all homomorphisms from \( G \) into the multiplicative group
    of roots of unity in \( \mathbb{C} \) with the group operation defined
    as
    \[
        \left( \varphi \psi \right)(g) = \varphi(g) \psi(g)
    \]
    for all \( \varphi, \psi \in \hat{G} \) and \( g \in G\), where the latter
    multiplication takes place in \( \mathbb{C} \).
\end{definition}

\begin{remark}
    We will often call the elements of the dual group
    \emph{characters}. If a group is
    isomorphic to its complement it will be called \emph{self-dual}.
\end{remark}

The following proposition will be used throughout the thesis, as almost
all of the groups we will be dealing with will be finite and abelian.

\begin{proposition}
    If \( G \) is finite and abelian, then it is self-dual.
\end{proposition}

\begin{proof}
    First we will prove the statement for finite cyclic groups. Let \( G =
    \left\langle g \right\rangle \), \( \operatorname{ord}(g) = n \). We will
    prove that the dual is cyclic of order \( n \). Let \( \chi\colon
        G \rightarrow
    \mathbb{C} \) be a character. Since the group is cyclic, the character is
    entirely determined by the value it takes on the generator \( g \). Since
    \( g^{n} = e \), obviously \( \chi(g)^{n} = 1 \), so \( \chi(g)^{n} \in
        \mu_{n} = \left\{ \exp\left( 2k \pi i / n \right) \colon k =
            0, 1, \dots, n-1
    \right\}\). For each possible choice of \( \chi(g) \in \mu_{n} \) we get a
    different character. It is then easy to see that \( \chi\colon G \rightarrow
    \mathbb{C} \) such that \( \chi(g) = \exp\left( 2 \pi i / n \right) \)
    generates this group, i.e. \( \hat{G} = \left\{ \chi, \chi^{2},
    \dots, \chi^{n-1}, \chi^{n} = e_{\hat{G}} \right\} \).

    Now since every finite abelian group can be written as a finite product
    of finite cycles, all we need to do is check whether the statement
    holds for a product of two groups \( A, B \) that are themselves
    self-dual, i.e. \( \widehat{A \times B} \cong \widehat{A} \times
    \widehat{B} \). This is self-evident as every character
    \( \chi \in \widehat{A \times B} \) is determined by its restrictions
    to \( A \times \{e_{B}\} \) and \( \left\{ e_{A} \right\} \times B \), since
    \[
        \chi(a, b) = \chi((a, e_{B})(e_{A}, b)) = \chi(a, e_{B})\chi(e_{A}, b).
    \]
    Thus \( \chi \leftrightarrow (\chi|_{A}, \chi|_{B})
    \) is an isomorphism.
\end{proof}

\section{\(C^{*}\)-algebras}

In the following section we will introduce some notions stemming from the
theory of \( C^{*} \)-algebras. We assume the Reader is familiar with
the notion of a \(C^{*}\)-algebra. For reference, see \textcite{blackadar}.

Multiplication in a \(C^{*}\)-algebra \( A \) is a bilinear map
\[
    m \colon A \times A \to A.
\]
By the universal property of the tensor product, \( m \) extends uniquely
to a linear map
\[
    m^{*} \colon A \otimes A \to A.
\]
Throughout this thesis, we will use the same notation \( (m) \) for both
the bilinear multiplication and its unique linear extension, relying
on the context to distinguish between the two.

\begin{definition}[Linear functional]
    A \emph{linear functional} on a \(C^{*}\)-algebra \( A \) is a
    linear function \( \psi\colon A \rightarrow \mathbb{C} \).
\end{definition}

\begin{remark}
    We will call linear functionals simply as \emph{functionals}.
\end{remark}

The following notions will be useful for defining an inner product on
\(C^{*}\)-algebras. Positiveness and faithfulness are all we need to
define a norm on a finite-dimensional \(C^{*}\)-algebra equipped with
a suitable functional.

\begin{definition}[Positive functional]
    A functional \( \psi \) on a \(C^{*}\)-algebra \( A \) is called
    \emph{positive}, if
    \[
        \psi(a^{*}a) \geq 0
    \]
    for any \( a \in A \).
\end{definition}

\begin{definition}[Faithful functional]
    A positive functional \( \psi \) on a \(C^{*}\)-algebra \( A \) is
    called \emph{faithful}, if
    \[
        \psi(a^{*}a) = 0 \Leftrightarrow a = 0
    \]
    for any \( a \in A \).
\end{definition}

% TODO: What literature?
\begin{remark}
    In the literature, the functional \( \psi \) is often called a \emph{state}.
    % For example see \textcite{gromada22}.
\end{remark}

The following notion is akin to the notion of a group action on a set.

\begin{definition}[Group action on a \(C^{*}\)-algebra]
    By \emph{group action} of a group \( G \) on a \(C^{*}\)-algebra \(
    A \) we mean a group homomorphism
    \[
        \alpha\colon G \rightarrow \operatorname{Aut}(A).
    \]
    We will write \( g \cdot a \) instead of \( \alpha(g)(a) \) then the action
    in question is clear.
\end{definition}

\subsection{Crossed-product \(C^{*}\)-algebra}

% TODO: rewrite this subsection

In this section we will introduce the notion of a crossed-product
\(C^{*}\)-algebra. We will do so in a limited manner, only to the
extent to which the theory of such products is necessary for our
purposes. For a more general reference we refer the Reader to
\textcite{blackadar}.

The following definitions are based off of \textcite{sierakowski09}.

\begin{definition}[\( C^{*} \)-dynamical system]
    A \( C^{*} \)-dynamical system is a triple \( (A, G, \alpha) \),
    where \( A \) is a \(C^{*}\)-algebra, \( G \) is a group, and \(
    \alpha\colon G \rightarrow \operatorname{Aut}(A) \) is a group
    action on \( A \).
    For a \( C^{*} \) dynamical system \( (A, G, \alpha) \) with \( G
    \) finite we define
    \[
        C_{c}(G, A, \alpha) = \left\{ f\colon G \rightarrow A \colon
        \operatorname{supp}(f) \text{ is finite}\right\},
    \]
    i.e. the space of all finitely supported functions on \( G \) with values
    in \( A \).
\end{definition}

\( C_{c}(G, A, \alpha) \) is equipped with both a product and a \( *
\)-operation in such a way that the action becomes "inner", i.e. \( g
\cdot a = u_{g} a u^{*}_{g} \) for all \( g \in G \), \( a \in A \).
More precisely we define those two operations in the following way:
\[
    ab = \sum_{g,h \in G}
\]

\begin{definition}[Covariant representation]
    Fix a dynamical system \( (A, G, \alpha) \). A \emph{covariant
    representation} \( (\pi, u, H) \) of \( (A, G, \alpha) \) consist of a
    unitary representation \( u\colon G \rightarrow B(\mathcal{H}) \)
    of \( G \) and
    a representation \( \pi\colon A \rightarrow B(\mathcal{H}) \) of
    \( A \) such
    that \( u(g) \pi(a) u(g)^{*} = \pi(g \cdot a) \) for every \( g \in G \)
    and \( a \in A \).
\end{definition}

\begin{definition}[Full \(C^{*}\)-algebra norm on \( C_{c}(G, A, \alpha) \)]
    The full \(C^{*}\)-algebra norm on \( C_{c}(G, A, \alpha) \) is given by
    \[
        \left\| \cdot \right\| = \sup \left\| (\pi \times u)(\cdot) \right\|,
    \]
    where the supremum is taken over all covariant representations \(
    (\pi, u, H) \)
    of \( (A, G, \alpha) \).
\end{definition}

\begin{definition}[Crossed product \(C^{*}\)-algebra]
    The crossed product of a \(C^{*}\)-algebra and a discrete group \( G \) is
    the completion of \( C_{c}(G, A, \alpha) \) with respect to the full norm.
    We denote the crossed product by \( \tilde{B}
    \rtimes_{\hat{\alpha}} \hat{G} \).
\end{definition}
